\origpage{24}{ROBERTSON AND SEYMOUR}
\noindent drawing of $H$ can be extended to a drawing of $G$. This is a much easier question. Every vertex of $H$ is surrounded by a ``large planar neighbourhood'' in $\Sigma$, and the remaining parts of $G$ attached in this neighbourhood have to be drawn within it, if the extension is possible; and so the problem essentially becomes a problem about extending drawings in discs. In fact, we can make the problem even easier, because we are only concerned with the case when $H$ is a subdivision of a simple 3-connected graph. The advantage of this hypothesis is that we get better, more useful, obstructions, because we cannot have ``made a mistake'' in the embedding of $H$. If this embedding of $H$ does not extend to one for all of $G$, we cannot locally redraw $H$ so that the embedding does extend, because there are no other ways to draw $H$ (this is a consequence of a theorem of Whitney, as we shall see); and this makes the obstructions healthier, less trivial.

On the other hand, one obstruction might be an edge $e$ of $G$ not in $H$, with ends $u, v \in V(H)$, not in the same region, but lying in distinct regions of $H$ which are incident with a common edge $f$. If we draw $e$ in the natural way crossing $f$, we get a graph drawn in $\Sigma$ with one crossing. In our application we can only make use of such a graph if the four ends of $e, f$ are not cut off from the main part of $H$ by a 3-separation. Crossings which are cut off by 3-separations are of no use to us, and we must modify the definition of ``extending the embedding of $H$'' so that these no longer are obstructions. For this reason, we shall only be concerned with the question, when can the embedding of $H$ be extended to one of $G$ ``up to 3-separations''? (We shall make this precise later.) Then all the obstructions (there are only two) turn out to be objects that we can make use of in our application.

In order to give the reader a better idea of what we are aiming for, we first state a simplified special case of the result, as follows.

\stmt
(1.1) \textit{Let $G$ be a 4-connected graph, and let $H$ be a subgraph of $G$ which is a subdivision of a simple 3-connected graph. Let $\Sigma$ be a connected surface without boundary, not a sphere, and let $H$ be drawn in $\Sigma$. Suppose that every closed curve in $\Sigma$ that meets $H$ at most 208 times is null-homotopic. For such a curve $\phi$, let $\operatorname{ins}(\phi)$ denote the union of the points of $\Sigma$ appearing in $\phi$ together with all discs in $\Sigma$ with boundary in this set. Then either}

\hspace*{12pt}(i) \textit{the drawing of $H$ can be extended to one of $G$, or}

\hspace*{12pt}(ii) \textit{there is a path of $G$ with distinct ends $u, v \in V(H)$ and with no other vertex or edge in $H$, such that there is no closed curve $\phi$ meeting $H$ at most 3 times with $u, v \in \operatorname{ins}(\phi)$, or}

\hspace*{12pt}(iii) \textit{there is a circuit $C$ of $H$ bounding an open disc $\Lambda \subseteq \Sigma$, and a point $x$ of $\Lambda$, and four vertices $a, b, c, d$ of $C$ in order, such that}
